% Methods equations for the SWOTxAI baselines ladder.
% Standalone: compile with pdflatex. To use in the paper, \input the body
% between the BEGIN BODY / END BODY markers. Written 2026-11-03 from
% scripts/baselines_eval.py; keep in step with docs/methods-notes.md section 3.
\documentclass[11pt]{article}
\usepackage[margin=1in]{geometry}
\usepackage{amsmath,amssymb,booktabs,array}
\usepackage[T1]{fontenc}
\title{Evaluation and baseline equations}
\author{SWOTxAI}
\date{3 November 2026}
\begin{document}
\maketitle

% ---------------------------------------------------------------- BEGIN BODY

\section{Notation}

One row is one radar grid cell $c$ at one SWOT pass time $t$. The target is
the 25-hour centred mean of the HFRNet surface current at that cell and time,
$\mathbf{y} = (y_u, y_v)$. Inputs at the row:
\begin{align*}
\mathbf{g} &= (g_u, g_v) && \text{SWOT L3 geostrophic velocity at the centre cell,}\\
\mathbf{W} &= (W_u, W_v) && \text{ERA5 10-m wind at the pass time,}\\
\mathbf{d} &= (d_u, d_v) && \text{DUACS L4 geostrophic velocity interpolated to the row,}\\
\mathbf{x} &\in \mathbb{R}^{72} && \text{the $3\times3$ stencil of eight SWOT, wind and MDT variables}\\
&&& \text{(81 columns on the calval orbit, where GOES SST is included),}\\
\tilde{\mathbf{x}} &&& \text{the stencil standardised by training-column mean and standard deviation,}\\
&&& \text{missing values set to zero after standardising.}
\end{align*}
$\bar{\mathbf{y}}$ is the mean of $\mathbf{y}$ over the training rows of the
region; $\bar{\mathbf{y}}(c)$ is the mean over the training rows at cell $c$.
Rows are split in time: training cycles first, a buffer of $B$ days after the
last training pass (30 days on the science orbit, 7 on calval), then the
held-out passes.

\section{Scoring}

For one region and one component, with $n$ held-out rows, prediction $\hat y_i$
and held-out mean $\langle y\rangle$,
\begin{equation}
R^2 = 1 - \frac{\sum_{i=1}^{n}(y_i - \hat y_i)^2}{\sum_{i=1}^{n}(y_i - \langle y\rangle)^2},
\qquad
\mathrm{RMSE} = \sqrt{\tfrac{1}{n}\sum_{i=1}^{n}(y_i - \hat y_i)^2}.
\label{eq:r2}
\end{equation}
The sums run over every held-out cell and pass of the region pooled together,
so the denominator contains the between-cell spread of the mean flow as well
as the pass-to-pass variability. Confidence intervals are a block bootstrap
over (cycle, pass): 1000 resamples of whole passes with replacement, 2.5 and
97.5 percentiles.

The skill score against the per-cell climatology (eq.~\ref{eq:clim}) removes
the between-cell credit:
\begin{equation}
S = 1 - \frac{\sum_i (y_i - \hat y_i)^2}{\sum_i \big(y_i - \bar y(c_i)\big)^2},
\label{eq:skill}
\end{equation}
which is zero for the climatology itself and positive only when a model
knows something the map of the mean flow does not.

\section{Baselines}

Every fitted coefficient below is obtained by ordinary least squares on the
training rows of the region, separately for $u$ and for $v$ (each component
has its own coefficients of the same form), and applied unchanged to the
held-out rows. Only the $u$ form is written.

\paragraph{Fixed rules, no parameters.}
\begin{align}
\hat u &= \bar y_u && \text{regional training mean,} \label{eq:trmean}\\
\hat u &= g_u && \text{SWOT geostrophy, slope one,} \label{eq:geos}\\
\hat u &= d_u && \text{DUACS geostrophy, slope one,} \label{eq:duacs}\\
\hat u &= g_u + \epsilon\,(W_u\cos\theta + W_v\sin\theta) && \text{Ekman rule, } \epsilon = 0.03,\ \theta = 45^\circ, \label{eq:ekman}\\
\hat v &= g_v + \epsilon\,(-W_u\sin\theta + W_v\cos\theta) && \text{(wind rotated } \theta \text{ to the right).} \nonumber
\end{align}

\paragraph{Climatology (the map).}
\begin{equation}
\hat u = \bar y_u(c), \qquad
\bar y_u(c) = \frac{1}{n_c}\sum_{i \in \text{train}(c)} y_{u,i},
\label{eq:clim}
\end{equation}
the training-period mean at the row's own cell; a held-out cell with no
training rows falls back to $\bar y_u$.

\paragraph{Linear rungs.} With $(a, b, c, d, e)$ fitted per component:
\begin{align}
\hat u &= a + b\,W_u + e\,W_v && \text{wind only,} \label{eq:wind}\\
\hat u &= \bar y_u(c) + a + b\,W_u + e\,W_v && \text{climatology + wind, fitted to } y_u - \bar y_u(c), \label{eq:climwind}\\
\hat u &= g_u + a + b\,W_u + e\,W_v && \text{SWOT (slope one) + wind, fitted to } y_u - g_u, \label{eq:geoswind}\\
\hat u &= a + b\,g_u && \text{SWOT calibrated,} \label{eq:geoscal}\\
\hat u &= a + b\,g_u + c\,g_v + d\,W_u + e\,W_v && \text{SWOT calibrated + wind,} \label{eq:geoscalwind}\\
\hat u &= a + b\,d_u && \text{DUACS calibrated,} \label{eq:duacscal}\\
\hat u &= a + b\,d_u + c\,d_v + d\,W_u + e\,W_v && \text{DUACS calibrated + wind.} \label{eq:duacscalwind}
\end{align}

\paragraph{Per-cell linear (partially pooled).}
The five-parameter model of eq.~\ref{eq:geoscalwind} fitted at each cell on
standardised inputs and shrunk toward the regional fit. With $A_c$ the
$n_c \times 5$ design matrix of cell $c$ (a column of ones and the four
standardised inputs), $\mathbf{y}_c$ its training targets, and
$\boldsymbol\beta_g$ the regional least-squares coefficients,
\begin{equation}
\boldsymbol\beta_c = \big(A_c^{\top}A_c + \lambda I\big)^{-1}
\big(A_c^{\top}\mathbf{y}_c + \lambda\,\boldsymbol\beta_g\big),
\qquad \lambda = 5,
\label{eq:celllinear}
\end{equation}
so the prior is worth five training rows; cells with fewer than five complete
training rows use $\boldsymbol\beta_g$. With about 170 rows per cell on the
calval orbit the local fit dominates; with about ten on the science orbit the
prior does.

\paragraph{Ridge on the full stencil.}
\begin{equation}
\hat u = a + \boldsymbol\beta^{\top}\tilde{\mathbf{x}},
\qquad
\boldsymbol\beta = \arg\min_{\boldsymbol\beta}
\sum_{i \in \text{train}} \big(y_{u,i} - a - \boldsymbol\beta^{\top}\tilde{\mathbf{x}}_i\big)^2
+ \alpha\,\|\boldsymbol\beta\|^2,
\qquad \alpha = 1.
\label{eq:ridge}
\end{equation}
This is the linear ceiling on the champion's own inputs.

\paragraph{AlphaEvolve champion.}
\begin{equation}
(\hat u, \hat v) = f_{\theta}(\tilde{\mathbf{x}}),
\label{eq:champion}
\end{equation}
a multilayer perceptron with joint $(u,v)$ output, trained by minimising the
squared error masked to finite targets, with early stopping on a validation
split; the architecture and training schedule are the candidate program
selected by the evolutionary search (\texttt{candidate.py} of the run).

\section{Alternative truth tiers}

The same predictions are also scored against three other versions of the
radar current built from the hourly record $u(t)$ at each cell: the raw hourly
value; and two zero-phase Butterworth low-passes of order 2, applied forward
and backward (\texttt{filtfilt}), with cutoff periods of 18 and 48 hours. The
headline 25-hour target is the centred boxcar
\begin{equation}
\bar u_{25}(t) = \frac{1}{25}\sum_{k=-12}^{12} u(t + k\,\text{h}),
\label{eq:boxcar}
\end{equation}
whose power response $|H(f)|^2$ has its first null at $f = 1/25\ \text{h}^{-1}
= 0.96$ cycles per day. Tiers are scored on the rows where every tier is
finite.

% ------------------------------------------------------------------ END BODY

\end{document}
